% =============================================================================
% COMPLÉMENTS CHAPITRE 1 — FICHIER DE RÉFÉRENCE
% =============================================================================
% Généré le 24 mai 2026
% Ce fichier contient TOUS les blocs à intégrer dans les sections §1.7 à §1.11
% Chaque bloc est précédé d'instructions d'intégration précises
% =============================================================================


% #############################################################################
%
%                              SECTION §1.7
%                    L'ÉLECTRIFICATION ET LE SUBSTRAT COMMUN
%
% #############################################################################
% Fichier cible : ch1_s07_electrification.tex
% =============================================================================


% -----------------------------------------------------------------------------
% BLOC §1.7-A : TRAJECTOIRES EUROPÉENNES
% -----------------------------------------------------------------------------
% POSITION : Insérer APRÈS la ligne 31 (fin du paragraphe sur Dennard/scaling)
%            AVANT le commentaire "% --- Les subsections continuent ici..."
% SUPPRIMER : Les lignes 33-43 (commentaires obsolètes sur le bilan à rédiger)
% -----------------------------------------------------------------------------

\subsection{Berlin, Londres et la divergence des trajectoires}\label{subsec:trajectoires_europe}

\dimmark{D2/0-GOV}{+}%
Le système Edison que la section précédente a décrit ne s'est pas déployé à l'identique dans les trois capitales où il fut transféré. \textcite{hughes_networks_1983} a documenté les trajectoires divergentes de New York, Berlin et Londres entre 1882 et 1914~; leur comparaison révèle que la forme institutionnelle, et non la supériorité technique, détermine le rythme d'électrification. À Berlin, Emil Rathenau, qui avait acquis les droits Edison pour l'Allemagne après avoir visité l'exposition de Paris en 1881, fonde la \emph{Deutsche Edison Gesellschaft} (1883), devenue \emph{Allgemeine Elektricitäts-Gesellschaft} (AEG) en 1887. La concession municipale de 1884 lui accorde le monopole de la distribution dans un périmètre défini, en échange du partage des revenus avec la ville~; AEG contrôle simultanément la fabrication des équipements et l'exploitation du réseau, une intégration verticale que ni Edison à New York ni aucun concurrent à Londres n'atteindra\footnote{Siemens \& Halske, le concurrent berlinois d'AEG, adopte une stratégie complémentaire plutôt que frontale~: les deux firmes se partagent le marché allemand et coopèrent dans les projets d'infrastructure à grande échelle. \textcite{hughes_networks_1983}, ch.~7, documente cette coordination comme un trait distinctif du modèle allemand.}. En 1914, la Berliner Elektrizitäts-Werke (BEW), filiale d'AEG, dessert un rayon de trente kilomètres avec six centrales et une capacité installée de 155\,000~kW~; le taux de charge (\emph{load factor}) y est supérieur à celui de Chicago et bien au-dessus de celui de Londres.

\dimmark{0-GOV}{+}%
À Londres, la trajectoire fut inverse. L'\emph{Electric Lighting Act} de 1882, promu par Joseph Chamberlain au Board of Trade, imposa aux compagnies privées une clause de rachat municipal à vingt et un ans au prix des actifs isolés (\emph{scrap iron clause})~: la municipalité pouvait racheter câbles, dynamos et bâtiments à leur valeur séparée, sans prime pour la valeur du système en fonctionnement\footnote{L'expression \emph{scrap iron} provient des débats parlementaires de 1882~: le secrétaire du Board of Trade expliqua que le rachat se ferait à la valeur des «~briques comme briques et du mortier comme mortier~» (\emph{bricks as bricks and mortar as mortar}). Voir \textcite{hughes_networks_1983}, p.~57.}. Cette clause gela l'investissement privé~: pourquoi immobiliser du capital dans une infrastructure dont la valeur serait amputée au moment du rachat~? En 1883, soixante-neuf autorisations provisoires furent accordées~; en 1884, quatre~; en 1885, aucune. La station Edison de Holborn Viaduct, ouverte en 1882, fut abandonnée par Edison \& Swan en 1886. Quand le Parlement amenda la loi en 1888 pour porter la durée de concession à quarante-deux ans, il était trop tard~: les villes avaient acquis le droit de gérer elles-mêmes l'éclairage, et la fragmentation s'enracina. En 1918, Londres comptait encore soixante-cinq compagnies d'électricité, dix fréquences différentes et vingt-quatre tensions~; l'unification ne viendrait qu'avec le \emph{Central Electricity Board} de 1926.

\dimmark{D2/D10}{*}%
La comparaison Hughes établit un résultat qui dépasse le cas de l'électricité. La trajectoire technique du réseau n'explique pas sa forme institutionnelle~: Berlin et Londres partaient du même artefact (le système Edison), disposaient d'ingénieurs compétents et de capitaux disponibles~; c'est la configuration politique qui discrimina les trajectoires. À Berlin, la coalition Rathenau--banques d'investissement--municipalité avait le pouvoir d'imposer une concession stable~; à Londres, la coalition aristocratie--science--industrie se heurta à des intérêts contradictoires au Parlement. Hughes résume~: \og~The most penetrating explanation for the failure in London and the success in Berlin is neither technological nor economic; it is political.\fg{}\footnote{\textcite{hughes_networks_1983}, p.~78.} Ce résultat est exactement parallèle à celui que la section~\ref{sec:telegraph} avait établi pour le télégraphe~: sept trajectoires nationales convergeant vers le monopole par sept mécanismes différents, dont la variété tient aux configurations politiques locales et non aux propriétés techniques du réseau. La dimension \textbf{0-GOV} n'est pas un résidu~; elle est, dans les deux cas, le facteur déterminant de la forme institutionnelle.

\emergence{La divergence Berlin/Londres confirme que les réseaux électriques, comme les réseaux télégraphiques, convergent vers des formes institutionnelles comparables (monopole régulé, intégration verticale, contrôle étatique) par des trajectoires politiquement déterminées. La supériorité technique n'explique ni le succès de Berlin ni l'échec de Londres~: c'est la capacité des coalitions d'acteurs à stabiliser un cadre institutionnel favorable à l'investissement long qui discrimine les trajectoires. Le mécanisme vaut pour 0-TER~: là où l'électrification est retardée par la fragmentation, les techniques de l'information restent plus longtemps sur substrat mécanique ou manuel, ce qui retarde aussi leur couplage au réseau de puissance.}


% -----------------------------------------------------------------------------
% BLOC §1.7-B : BILAN — SÉQUENCE D'ÉLECTRIFICATION T→C→S
% -----------------------------------------------------------------------------
% POSITION : Insérer APRÈS la fin de \subsection{Enigma...} (ligne ~97)
%            AVANT le commentaire de fin "% ===" (ligne ~99)
% -----------------------------------------------------------------------------

\subsection{Bilan~: la séquence d'électrification des trois fonctions}\label{subsec:bilan_electrification}

\dimmark{0-TER}{*}%
Le parcours de cette section a fait apparaître une séquence qui n'avait pas été nommée jusque-là. Les trois fonctions informationnelles~--- transmettre, stocker, calculer~--- ne s'électrifient pas au même moment~:

\begin{itemize}
\item \textbf{TRANSMIT} s'électrifie dès l'origine~: le télégraphe électrique (1837) \emph{est} électrique par définition. La pile voltaïque, puis le réseau à courant continu et alternatif, portent le signal dès le premier jour. L'électrification de la transmission n'est pas une transition~; c'est une condition initiale.

\item \textbf{COMPUTE} reste mécanique ou manuel pendant plus d'un siècle. La règle à calcul, la calculatrice de Thomas de Colmar, les tabulatrices Hollerith ne tirent pas leur énergie du réseau électrique~: elles sont actionnées à la main ou par des moteurs électriques dont la consommation est marginale dans le bilan d'exploitation. L'électronique de puissance n'entre dans le calcul qu'avec l'ENIAC (1946), soit cent neuf ans après Morse.

\item \textbf{STORE} bascule en dernier. Le papier, le carton perforé, la cire du phonographe, le film photographique, le disque de gomme-laque stockent l'information sans énergie~: l'inscription est une déformation physique permanente du support, et la lecture en est la traduction mécanique ou optique. La mémoire DRAM (Intel 1103, 1970) est la première technologie de stockage qui \emph{exige} un flux d'énergie continu~--- plusieurs centaines de rafraîchissements par seconde~--- pour maintenir son contenu. Avant elle, stocker était gratuit entre l'écriture et la lecture~; après elle, stocker consomme en permanence.
\end{itemize}

\dimmark{0-TER}{*}%
Cette séquence~--- T~(1837) $\rightarrow$ C~(1946) $\rightarrow$ S~(1970)~--- a une conséquence directe pour l'argument de la thèse. Tant que STORE reste passif, le coût énergétique de l'information est ponctuel~: il se produit à la transmission et au calcul, pas entre les deux. L'information peut dormir sur une étagère sans consommer de watt. Le régime énergétique est additif~: chaque usage ajoute une dépense discrète, mais il n'y a pas de fond continu. C'est ce régime que les sections §\ref{sec:telegraph} à §\ref{sec:wireless} ont traversé sans avoir besoin de le nommer, parce qu'il ne posait pas de problème~: l'énergie de l'information était invisible parce qu'elle était négligeable.

\dimmark{0-TER}{+}%
Le basculement vers un régime permanent s'opère en deux temps. Le premier temps est la DRAM (1970)~: la mémoire de travail des ordinateurs devient un flux continu de rafraîchissement. Le second temps est le serveur web (1991) et le datacenter (années~2000)~: l'exigence de disponibilité permanente (\emph{always-on}) transforme STORE d'une ressource ponctuelle en une infrastructure continue. Un fichier stocké sur un serveur web n'existe que si le serveur est allumé~; une page web accessible à tout moment exige une alimentation à tout moment. La section~\ref{sec:convergence} documentera ce basculement~; la section~\ref{sec:contemporain} en mesurera les conséquences. Ce que la présente section établit est le point de départ~: jusqu'en 1970, le substrat électrique des techniques de l'information est une condition d'existence, pas un coût d'exploitation continu. L'occultation de la dimension 0-TER dans les cas du XIX\textsuperscript{e}~siècle n'est pas une erreur d'observation~: elle reflète un régime physique où l'énergie de l'information était réellement négligeable.

\questionch{3}{Comment quantifier la consommation d'énergie de chaque fonction (TRANSMIT, STORE, COMPUTE) à travers les périodes identifiées par ce chapitre~? Le croisement des séries de capacités informationnelles de \textcite{hilbert_worlds_2011} avec les séries d'efficacité énergétique de \textcite{koomey_implications_2011} permettrait-il de reconstituer la consommation agrégée par fonction et par décennie~? Ce croisement n'a jamais été publié. Le chapitre~3 de cette thèse le tentera.}


% -----------------------------------------------------------------------------
% BLOC §1.7-C : TABLEAU RÉCAPITULATIF DES DIMENSIONS
% -----------------------------------------------------------------------------
% POSITION : Insérer IMMÉDIATEMENT APRÈS le bloc §1.7-B (bilan)
%            AVANT le commentaire de fin "% ==="
% -----------------------------------------------------------------------------

\begin{table}[htbp]
\caption{Dimensions économiques de l'électrification (§\ref{sec:electrification}).}\label{tab:electrification-dimensions}
\centering\small
\renewcommand{\arraystretch}{1.15}
\begin{tabularx}{\textwidth}{@{}cXc@{}}
\toprule
\textbf{Dim.} & \textbf{Ce que §\ref{sec:electrification} établit} & \textbf{Niveau} \\
\midrule
D2 & Monopole naturel par rendements d'échelle (Insull)~; fragmentation = retard (Londres vs Berlin). & ★ \\
D3 & Capital de réseau irréversible (centrale + distribution)~; intégration verticale AEG. & + \\
D4 & Création d'emplois dans la maintenance et l'exploitation~; déplacement du travail vers la supervision (Nautical Almanac). & + \\
D5 & Demande dérivée (éclairage, puis moteur, puis information)~; diversification des usages pour optimiser le \emph{load factor}. & ★ \\
D6 & Le compteur kilowattheure transforme un flux continu en signal discret facturable~: l'information rend le système de puissance viable. & + \\
D7 & Standardisation des fréquences et tensions comme condition d'interconnexion~; fragmentation londonienne (24 tensions, 10 fréquences). & + \\
D8 & Cycle investissement lourd $\rightarrow$ consolidation (Chicago Edison $\rightarrow$ Commonwealth Edison)~; financement par les banques d'investissement (Deutsche Bank, Drexel Morgan). & ★ \\
D10 & Lock-in DC vs AC (guerre des courants)~; \emph{technological momentum} de Hughes~; \emph{scaling} de Dennard comme bouclier d'occultation. & ★ \\
\midrule
0-TER & Séquence T\,(1837) $\rightarrow$ C\,(1946) $\rightarrow$ S\,(1970)~; régime passif du stockage pré-DRAM~; STORE permanent = consommation continue. & ★ \\
0-GOV & Régulation des monopoles naturels (\emph{public utility commissions})~; divergence Berlin/Londres par la configuration politique, non technique. & ★ \\
\bottomrule
\end{tabularx}
\renewcommand{\arraystretch}{1.0}
\end{table}



% #############################################################################
%
%                              SECTION §1.8
%                 LE CALCUL PROGRAMMABLE ET LA RÉVOLUTION ÉLECTRONIQUE
%
% #############################################################################
% Fichier cible : ch1_s08_calcul_programmable.tex
% =============================================================================


% -----------------------------------------------------------------------------
% BLOC §1.8-A : TRANSITION TUBE-TRANSISTOR-IC
% -----------------------------------------------------------------------------
% POSITION : Insérer APRÈS la ligne 30 (fin de l'introduction)
%            AVANT \subsection{La demande de calcul...} (ligne 33)
% -----------------------------------------------------------------------------

% -----------------------------------------------------------------------------
\subsection{Du tube au transistor~: la transition matérielle (1947--1965)}\label{subsec:transistor}
% -----------------------------------------------------------------------------

\dimmark{D10}{+}%
Le 23~décembre 1947, aux Bell Telephone Laboratories, John Bardeen, Walter Brattain et William Shockley réalisent le premier transistor à pointe de contact. Le dispositif, encore fragile et coûteux, remplace le tube à vide dans sa fonction d'amplification et de commutation~: plus petit (quelques millimètres contre plusieurs centimètres), plus fiable (pas de filament à chauffer), et surtout beaucoup moins gourmand en énergie (milliwatts contre watts). Le brevet est déposé en 1948~; le prix Nobel couronne les trois inventeurs en 1956. La trajectoire qui va du premier transistor au circuit intégré couvre moins de quinze ans, une compression temporelle sans équivalent dans les cas précédents.

\dimmark{D5/D8}{+}%
Cette compression n'est pas le fait du marché civil. \textcite{flamm_creating_1988} documente que la demande militaire absorbe la quasi-totalité de la production de transistors dans les premières années. Le programme de miniaturisation des missiles balistiques intercontinentaux (Minuteman) exige des composants légers et fiables~; la Navy et l'Air Force achètent 100\,\% de la production américaine de circuits intégrés en 1962. La courbe d'apprentissage qui en résulte fait chuter le prix du transistor de 45~dollars en 1954 à 2~dollars en 1960~--- une division par vingt en six ans. Le marché civil hérite d'une technologie dont le coût de développement a été socialisé par le budget de la défense.

\dimmark{D10}{+}%
Le circuit intégré (1958--1965) franchit une étape supplémentaire. Jack Kilby chez Texas Instruments et Robert Noyce chez Fairchild Semiconductor déposent des brevets indépendants la même année (1958-1959) pour l'intégration de plusieurs transistors sur un même substrat de silicium. La querelle de priorité qui s'ensuit~--- tranchée par un accord de licences croisées en 1966~--- importe moins que la conséquence architecturale~: le circuit intégré permet de graver des milliers, puis des millions de transistors sur une puce unique, réduisant les connexions externes et les points de défaillance. Noyce fonde Intel en 1968 avec Gordon Moore~; le premier microprocesseur commercial (Intel 4004) sort en 1971. La trajectoire de la loi de Moore (§\,\ref{subsec:moore_law}) commence ici.

\dimmark{0-TER}{+}%
Le profil énergétique de cette transition est paradoxal. L'efficacité par opération s'améliore de façon spectaculaire~: le ratio joules/opération chute de six ordres de grandeur entre l'ENIAC (1946) et l'IBM~360 (1964)\footnote{L'ENIAC consommait environ 150\,kW pour 5\,000 additions par seconde, soit 30~joules par opération. L'IBM~360/50 (1964) consommait environ 20\,kW pour 500\,000 opérations par seconde, soit 0,04~joule par opération. Le gain est de l'ordre de $10^{3}$. Entre l'IBM~360 et l'Intel~4004 (1971), le gain est encore de trois ordres de grandeur. Sources~: \textcite{nordhaus_two_2007}, \textcite{koomey_implications_2011}.}. Mais la puissance totale installée augmente~: les mainframes des années 1960 consomment 20 à 50~kW chacun, et leur nombre croît. Le paradoxe est déjà en place~: l'efficacité par unité s'améliore, la consommation agrégée croît parce que le volume de calcul croît plus vite que l'efficacité. C'est le mécanisme que \textcite{jevons_coal_1865} avait identifié pour le charbon et que la littérature contemporaine nomme \emph{rebound effect}~; il s'appliquera au calcul pendant les cinq décennies suivantes.


% -----------------------------------------------------------------------------
% BLOC §1.8-B : CONSENT DECREE ET UNBUNDLING
% -----------------------------------------------------------------------------
% POSITION : Insérer APRÈS la fin de \subsection{La loi de Moore...} (ligne ~89)
%            La section moore_law semble tronquée ; ce bloc vient après
% -----------------------------------------------------------------------------

% -----------------------------------------------------------------------------
\subsection{Le \emph{consent decree} et la naissance du logiciel comme marchandise}\label{subsec:consent_decree}
% -----------------------------------------------------------------------------

\dimmark{D2/0-GOV}{*}%
Le 14~janvier 1949, le Department of Justice dépose une plainte antitrust contre AT\&T, accusant le monopole téléphonique d'avoir utilisé sa position dominante pour exclure les concurrents du marché des équipements. L'affaire se conclut le 24~janvier 1956 par un \emph{consent decree} dont les conséquences dépassent largement le secteur des télécommunications. AT\&T conserve son monopole sur le service téléphonique, mais doit licencier l'ensemble de ses brevets~--- y compris le transistor~--- à tout demandeur, moyennant des redevances raisonnables. Texas Instruments, Fairchild Semiconductor et la future Intel bénéficient directement de cette diffusion forcée~: la technologie du transistor, développée aux frais des abonnés téléphoniques, devient accessible à l'industrie électronique tout entière\footnote{\textcite{temin_fall_1987} documente l'histoire du \emph{consent decree} de 1956 et ses effets sur la structure de l'industrie américaine des télécommunications. La clause de licence obligatoire est l'un des mécanismes par lesquels l'antitrust américain a involontairement accéléré la diffusion des technologies de l'information.}.

\dimmark{D2/D3}{*}%
Treize ans plus tard, le mécanisme se reproduit pour IBM. En janvier 1969, le dernier jour de l'administration Johnson, le Department of Justice ouvre une enquête antitrust contre IBM, accusant la firme de monopoliser le marché des ordinateurs en liant (\emph{bundling}) matériel, logiciel et services dans une offre indissociable. Le 23~juin 1969, six mois après l'ouverture de l'enquête et avant même toute décision judiciaire, IBM annonce l'\emph{unbundling}~: désormais, le matériel, le logiciel et les services seront facturés séparément. La décision est présentée comme volontaire~; elle anticipe une contrainte judiciaire que la firme préfère devancer. L'affaire elle-même traînera jusqu'en 1982, date à laquelle l'administration Reagan l'abandonnera~; mais le mal est fait, du point de vue d'IBM. Le logiciel est devenu une marchandise autonome.

\dimmark{D3}{*}%
Cette séparation a une conséquence comptable qui prépare la section~\ref{sec:convergence}. Avant 1969, le logiciel n'existe pas comme catégorie économique distincte~: il est inclus dans le prix de la machine, sans valeur comptable propre. Après 1969, le logiciel devient un bien séparé, appropriable, valorisable. Les premières entreprises de logiciel indépendant~--- Computer Associates (1976), SAP (1972), Oracle (1977)~--- naissent dans le sillage de l'\emph{unbundling}. \textcite{campbell-kelly_software_2003} documente cette naissance comme la création d'une industrie qui n'existait pas dix ans plus tôt. Le capital intangible de \textcite{corrado_intangible_2005}, dont la mesure structure la comptabilité nationale depuis les années 2000, trouve son origine institutionnelle dans une décision antitrust de 1969.

\dimmark{D3/0-TER}{+}%
Le point analytique pour la thèse est le suivant. L'\emph{unbundling} crée une asymétrie entre la visibilité comptable du logiciel et l'invisibilité de son substrat énergétique. Le logiciel, une fois séparé du matériel, peut être valorisé comme actif immatériel~: il figure dans les bilans, il est amorti, il génère des revenus de licence. Mais le logiciel n'existe que s'il est exécuté, et l'exécution consomme de l'électricité. Cette consommation n'est pas attachée au logiciel dans les comptes~: elle est noyée dans les charges d'exploitation de l'utilisateur, pas dans le coût de production du vendeur. L'\emph{unbundling} de 1969 inaugure ainsi une dissociation comptable entre la valeur du logiciel (visible, appropriable) et le coût énergétique de son usage (diffus, non imputable), dissociation qui structure encore aujourd'hui la comptabilité du secteur numérique.


% -----------------------------------------------------------------------------
% BLOC §1.8-C : TABLEAU EFFICACITÉ ÉNERGÉTIQUE (données chiffrées)
% -----------------------------------------------------------------------------
% POSITION : À intégrer dans le bilan (BLOC §1.8-D), après "Deuxième rupture"
% -----------------------------------------------------------------------------

\begin{table}[htbp]
\caption{Efficacité énergétique du calcul~: jalons 1946--1971.}\label{tab:efficacite-calcul}
\centering\small
\renewcommand{\arraystretch}{1.15}
\begin{tabular}{@{}llrrr@{}}
\toprule
\textbf{Année} & \textbf{Machine} & \textbf{Puissance (kW)} & \textbf{Opérations/s} & \textbf{Ops/kWh} \\
\midrule
1946 & ENIAC & 150 & 5\,000 & $1{,}2 \times 10^{2}$ \\
1954 & IBM 704 & 30 & 40\,000 & $4{,}8 \times 10^{3}$ \\
1964 & IBM 360/50 & 20 & 500\,000 & $9{,}0 \times 10^{4}$ \\
1971 & Intel 4004 & 0,0005 & 60\,000 & $4{,}3 \times 10^{11}$ \\
\bottomrule
\end{tabular}
\renewcommand{\arraystretch}{1.0}
\par\smallskip
\footnotesize Sources~: \textcite{nordhaus_two_2007}, \textcite{koomey_implications_2011}. Ops/kWh = opérations par kilowattheure.
\end{table}


% -----------------------------------------------------------------------------
% BLOC §1.8-D : BILAN — TROIS RUPTURES ET UN PARADOXE
% -----------------------------------------------------------------------------
% POSITION : Insérer APRÈS le bloc §1.8-B (consent decree)
% -----------------------------------------------------------------------------

% -----------------------------------------------------------------------------
\subsection{Bilan~: trois ruptures et un paradoxe}\label{subsec:bilan_computing}
% -----------------------------------------------------------------------------

\dimmark{0-TER}{*}%
La période 1945--1975 voit se produire trois ruptures simultanées dont la conjonction prépare les développements des sections suivantes.

\paragraph{Première rupture~: l'architecture à programme enregistré.} Le \emph{stored-program concept}, formalisé par von Neumann en 1945 et matérialisé dans l'EDVAC puis l'IAS Machine, unifie l'espace mémoire du programme et l'espace mémoire des données. Le programme devient donnée~: il peut être copié, modifié, transmis comme n'importe quel fichier. Cette fusion est la condition technique de l'industrie du logiciel~: sans elle, le logiciel resterait câblé dans la machine et ne pourrait pas être vendu séparément. Mais elle a aussi une conséquence énergétique~: la mémoire qui contient le programme consomme de l'énergie en permanence (rafraîchissement DRAM). Stocker un programme n'est plus gratuit énergétiquement, contrairement au carton perforé ou au câblage fixe.

\paragraph{Deuxième rupture~: la transition tube-transistor-IC.} L'efficacité énergétique par opération s'améliore de neuf ordres de grandeur entre l'ENIAC (1946) et l'Intel 4004 (1971). Cette amélioration résulte de la miniaturisation (moins de matière à commuter, moins d'énergie dissipée) et du passage à l'électronique à semi-conducteurs (pas de filament à chauffer). Le tableau~\ref{tab:efficacite-calcul} documente les jalons de cette trajectoire.

\paragraph{Troisième rupture~: l'\emph{unbundling} et la naissance du capital intangible.} L'\emph{unbundling} d'IBM (1969) crée le logiciel comme marchandise séparée du matériel. Cette séparation institutionnelle prépare la mesure du capital intangible de \textcite{corrado_intangible_2005}, qui structure aujourd'hui la comptabilité nationale. Elle prépare aussi la dissociation comptable entre la valeur du logiciel (visible, appropriable) et le coût énergétique de son usage (diffus, non imputable).

\dimmark{0-TER}{*}%
\paragraph{Le paradoxe.} Ces trois ruptures ont un effet combiné paradoxal. L'efficacité par opération s'améliore spectaculairement~; la consommation agrégée croît parce que le volume de calcul croît plus vite que l'efficacité. Le même mécanisme de rebond (\emph{rebound effect}) que \textcite{jevons_coal_1865} avait identifié pour le charbon s'applique au calcul~: chaque gain d'efficacité libère des ressources qui sont réinvesties en volume, de sorte que la consommation totale augmente. La section~\ref{sec:convergence} documentera comment ce paradoxe se déploie entre 1975 et 2007~; la section~\ref{sec:contemporain} montrera qu'il s'intensifie après 2007 avec l'IA générative.

\emergence{La période 1945--1975 voit trois ruptures simultanées~: l'architecture à programme enregistré qui fusionne STORE et COMPUTE~; la transition tube-transistor-IC qui améliore l'efficacité énergétique de neuf ordres de grandeur~; l'\emph{unbundling} qui crée le logiciel comme capital intangible séparé du matériel. Ces ruptures préparent la convergence numérique de §\,\ref{sec:convergence} et l'explosion de §\,\ref{sec:contemporain}. Le paradoxe qui les traverse~--- amélioration de l'efficacité unitaire, croissance de la consommation agrégée~--- est le mécanisme central de l'occultation de la dimension 0-TER dans les décennies qui suivent.}

\questionch{3}{Comment quantifier le rebond énergétique du calcul sur la période 1945--2020~? Le croisement des séries d'efficacité de \textcite{koomey_implications_2011} avec les séries de capacité de \textcite{hilbert_worlds_2011} et les statistiques de consommation électrique des datacenters de \textcite{masanet_recalibrating_2020} permettrait de tester l'hypothèse Jevons sur le calcul. Ce croisement n'a jamais été publié sous cette forme.}


% -----------------------------------------------------------------------------
% BLOC §1.8-E : TABLEAU RÉCAPITULATIF DES DIMENSIONS
% -----------------------------------------------------------------------------
% POSITION : Insérer APRÈS le bloc §1.8-D (bilan), AVANT "% ===" de fin
% -----------------------------------------------------------------------------

\begin{table}[htbp]
\caption{Dimensions économiques du calcul programmable (§\ref{sec:computing}).}\label{tab:computing-dimensions}
\centering\small
\renewcommand{\arraystretch}{1.15}
\begin{tabularx}{\textwidth}{@{}cXc@{}}
\toprule
\textbf{Dim.} & \textbf{Ce que §\ref{sec:computing} établit} & \textbf{Niveau} \\
\midrule
D2 & \emph{Consent decree} 1956 (AT\&T) et \emph{unbundling} 1969 (IBM)~: l'antitrust comme mécanisme de diffusion technologique. & ★ \\
D3 & \emph{Unbundling} 1969~: naissance du logiciel comme capital intangible séparé du matériel. & ★ \\
D5 & Demande militaire (tables de tir, cryptanalyse, simulation nucléaire, SAGE)~: l'État crée le marché que le privé ne pourrait financer. & ★ \\
D8 & L'État finance 2/3 de la R\&D dans les années 1950~; le marché civil prend le relais dans les années 1970--1980. & ★ \\
D10 & Loi de Moore comme dispositif de coordination performatif~; \emph{scaling} de Dennard comme bouclier d'occultation énergétique. Architectures alternatives bloquées (OGAS, Cybersyn). & ★ \\
\midrule
0-TER & Efficacité $\times 10^{9}$, mais consommation agrégée croît (rebond Jevons). Architecture von Neumann couple STORE et énergie (rafraîchissement DRAM). ENIAC = 150\,kW~; SAGE = 3\,MW par centre. & ★ \\
0-GOV & Antitrust (consent decree, unbundling)~; financement fédéral massif~; trajectoires alternatives bloquées par la configuration politique (OGAS, Cybersyn). & ★ \\
\bottomrule
\end{tabularx}
\renewcommand{\arraystretch}{1.0}
\end{table}



% #############################################################################
%
%                              SECTION §1.9
%                       LA CONVERGENCE SUR LE RÉSEAU
%
% #############################################################################
% Fichier cible : ch1_s09_convergence.tex
% =============================================================================


% -----------------------------------------------------------------------------
% BLOC §1.9-A : RENFORCEMENT DU PARAGRAPHE PARADOXE INTANGIBLE/TANGIBLE
% -----------------------------------------------------------------------------
% POSITION : Insérer APRÈS la ligne 99 (fin du paragraphe existant sur le paradoxe)
%            AVANT \questionch{} (ligne 101)
% Le paragraphe existant (lignes 98-99) est bon ; on ajoute juste ce complément
% -----------------------------------------------------------------------------

\dimmark{D3/0-TER}{*}%
Ce découplage comptable a une conséquence pratique~: les entreprises qui investissent massivement dans le capital intangible (logiciels, données, algorithmes) externalisent la consommation énergétique correspondante vers les fournisseurs de cloud et les opérateurs de datacenters, dont les coûts apparaissent comme des charges d'exploitation et non comme l'amortissement d'un actif. La structure comptable rend donc structurellement difficile l'imputation du coût énergétique de l'information à ses bénéficiaires économiques. Le chapitre~2 reprendra ce mécanisme sous la dimension D3~; le chapitre~3 tentera d'en estimer l'ampleur.


% -----------------------------------------------------------------------------
% BLOC §1.9-B : TABLEAU RÉCAPITULATIF DES DIMENSIONS
% -----------------------------------------------------------------------------
% POSITION : Insérer APRÈS la ligne 138 (fin du fichier, après \questionch{})
% -----------------------------------------------------------------------------

% -----------------------------------------------------------------------------
% Tableau récapitulatif des dimensions
% -----------------------------------------------------------------------------

\begin{table}[htbp]
\caption{Dimensions économiques de la convergence (§\ref{sec:convergence}).}\label{tab:convergence-dimensions}
\centering\small
\renewcommand{\arraystretch}{1.15}
\begin{tabularx}{\textwidth}{@{}cXc@{}}
\toprule
\textbf{Dim.} & \textbf{Ce que §\ref{sec:convergence} établit} & \textbf{Niveau} \\
\midrule
D1 & L'information numérique est non rivale~; le circuit qui la transporte est rival. Le paradoxe de la productivité (Solow 1987, David 1990) reflète un décalage de substrat, pas seulement d'organisation. & + \\
D2 & Plateformes \emph{winner-take-all} (Microsoft, Google)~; architecture ouverte IBM-PC $\rightarrow$ clonage $\rightarrow$ concentration sur l'OS et les applications. & + \\
D3 & Capital intangible $\geq$ tangible dès 1995 (\textcite{corrado_intangible_2009})~; \emph{unbundling} IBM (1969) comme origine institutionnelle~; découplage comptable entre valeur du logiciel et coût énergétique de son usage. & ★ \\
D5 & Paradoxe de la productivité (Solow 1987, David 1990)~; le gain de productivité n'apparaît qu'après recomposition organisationnelle complète. & ★ \\
D6 & Le Web (1993) crée un signal de prix instantané à l'échelle mondiale~; le moteur de recherche (Google, 1998) devient l'interface d'accès à l'information. & + \\
D7 & Standards ouverts (TCP/IP 1983, HTTP 1993) comme condition de la convergence~; modularité au sens de \textcite{langlois_modularity_2003}. & ★ \\
D8 & Cycle Perez~: dotcom = railway mania~; l'infrastructure (câbles, serveurs) survit au krach et sert la phase de déploiement. Secteur financier comme premier déployeur massif (1970--1990). & ★ \\
D10 & Open source (GNU 1984, Linux 1991) vs licensing massif (Microsoft)~: deux trajectoires coexistantes~; culture \emph{Whole Earth Catalog} comme trace structurelle. Bifurcations écartées (refus AT\&T 1971, échec Inouye 1994). & + \\
\midrule
0-TER & STORE devient permanent (serveur web \emph{always-on})~; premier régime de consommation continue~; la numérisation du réseau (1983--2000) est un double coût d'infrastructure, pas une substitution. & ★ \\
0-GOV & Privatisation NSFNET (1995)~; \emph{Framework for Global Electronic Commerce} (Clinton 1997) consacre l'autorégulation~; trajectoires alternatives (Inouye 1994) écartées. & + \\
\bottomrule
\end{tabularx}
\renewcommand{\arraystretch}{1.0}
\end{table}



% #############################################################################
%
%                              SECTION §1.10
%                       L'INFORMATION QUI CONSOMME
%
% #############################################################################
% Fichier cible : ch1_s10_information_consomme.tex
% =============================================================================


% -----------------------------------------------------------------------------
% BLOC §1.10-A : RÉSOLUTION DU \wip{} (ligne 52)
% -----------------------------------------------------------------------------
% POSITION : Ligne 52 — REMPLACER le \wip{} existant par cette footnote
% ANCIEN : \wip{Références à compléter~: ...}
% NOUVEAU : (footnote ci-dessous, à rattacher à la phrase précédente)
% -----------------------------------------------------------------------------

% REMPLACER LIGNE 52 PAR :
\footnote{Pour le récit géopolitique de la concentration du secteur, voir \textcite{miller_chip_2022}. Les données techniques sur les machines EUV proviennent des rapports annuels ASML (2020--2024) et des présentations investisseurs. Sur le financement consortial SEMATECH et ses effets sur la structure de l'industrie américaine, voir \textcite{irwin_learning_1994} et \textcite{flamm_targeting_1987}.}


% -----------------------------------------------------------------------------
% BLOC §1.10-B : DÉCISIONS SUR LES \dimmark{?}
% -----------------------------------------------------------------------------
% MODIFICATIONS PONCTUELLES À EFFECTUER :
%
% Ligne 93  : remplacer \dimmark{D5/D10}{?}  par \dimmark{D5/D10}{+}
% Ligne 98  : remplacer \dimmark{D5/0-TER}{?} par \dimmark{D5/0-TER}{+}
% Ligne 113 : remplacer \dimmark{D2/0-TER}{?} par \dimmark{0-TER}{+}
% Ligne 116 : remplacer \dimmark{0-TER}{?}    par \dimmark{0-TER}{+}
% Ligne 122 : remplacer \dimmark{0-TER}{?}    par \dimmark{0-TER}{+}
% Ligne 125 : remplacer \dimmark{0-TER}{?}    par \dimmark{0-TER}{+}
% Ligne 152 : remplacer \dimmark{D10}{?}      par \dimmark{D10}{+}
% -----------------------------------------------------------------------------


% -----------------------------------------------------------------------------
% BLOC §1.10-C : TABLEAU RÉCAPITULATIF DES DIMENSIONS
% -----------------------------------------------------------------------------
% POSITION : Insérer APRÈS la ligne 167 (fin de \subsection{Landauer})
%            AVANT le commentaire de fin "% ===" (ligne 169)
% -----------------------------------------------------------------------------

% -----------------------------------------------------------------------------
% Tableau récapitulatif des dimensions
% -----------------------------------------------------------------------------

\begin{table}[htbp]
\caption{Dimensions économiques de la période contemporaine (§\ref{sec:contemporain}).}\label{tab:contemporain-dimensions}
\centering\small
\renewcommand{\arraystretch}{1.15}
\begin{tabularx}{\textwidth}{@{}cXc@{}}
\toprule
\textbf{Dim.} & \textbf{Ce que §\ref{sec:contemporain} établit} & \textbf{Niveau} \\
\midrule
D2 & Concentration cloud (AWS, Azure, GCP)~; monopole ASML sur la lithographie EUV~; modèle Insull reproduit à l'échelle mondiale. & ★ \\
D3 & Externalisation du compute vers le cloud~: capital fixe $\rightarrow$ coût variable~; découplage comptable entre valeur logicielle et coût énergétique d'usage. & + \\
D4 & Déplacement du travail vers l'orchestration et la supervision (IA agentique)~; division internationale du travail semi-conducteur (design US, fab Taïwan, équipement NL). & + \\
D5 & Surplus comportemental (\textcite{zuboff_surveillance_2019})~; rebond potentiel de l'IA agentique (tâche $\rightarrow$ séquence d'inférences)~; scaling laws = amélioration $\propto$ compute. & ★ \\
D6 & AdWords/AdSense (2002--2003)~: marché de l'attention opérationnalisé~; requête ChatGPT $\approx$ 10$\times$ énergie recherche Google. & + \\
D8 & Contrats nucléaires 15--20 ans (Microsoft/Three Mile Island, Amazon/Susquehanna, Meta/Clinton)~: horizons d'investissement ICT et énergie synchronisés. Capex 2025 $>$ 320~Mrd\$. & ★ \\
D9 & Division géographique du travail semi-conducteur (US/TW/KR/NL/DE)~; restrictions export EUV vers Chine~; CHIPS Act, European Chips Act, Made in China 2025. & ★ \\
D10 & Loi de Moore comme dispositif de coordination (ITRS/IRDS)~; scaling laws IA comme nouvelle feuille de route~; niches alternatives (Fairphone, Te Hiku, DECODE). & ★ \\
\midrule
0-TER & Rupture Dennard ($\sim$2005)~; datacenters 415~TWh (2024) $\rightarrow$ 945~TWh (2030)~; TDP GPU 700~W (2024) $\rightarrow$ 1\,400~W (2026)~; EUV $\approx$ 1~MW/machine~; TSMC = 6~\% élec. Taïwan~; Bitcoin $\approx$ 150--200~TWh/an~; Landauer = plancher thermodynamique. & ★ \\
0-GOV & Fenêtre 2008 refermée sans bifurcation~; trois cadres régulatoires (\textcite{bradford_digital_2023})~: US \emph{market-driven}, EU \emph{rights-driven}, CN \emph{state-led}~; critique des \emph{twin transitions}. & ★ \\
\bottomrule
\end{tabularx}
\renewcommand{\arraystretch}{1.0}
\end{table}



% #############################################################################
%
%                              SECTION §1.11
%                       DES CAS AUX ARCHÉTYPES
%
% #############################################################################
% Fichier cible : ch1_s11_archetypes.tex
% =============================================================================


% -----------------------------------------------------------------------------
% BLOC §1.11-A : TABLEAU SYNOPTIQUE DIMENSIONS × PÉRIODES
% -----------------------------------------------------------------------------
% POSITION : Insérer APRÈS la ligne 83 (fin du tableau tab:ch1-questions)
%            AVANT la ligne 85 ("Le chapitre a révélé une structure...")
% -----------------------------------------------------------------------------

\begin{table}[p]
\caption{Matrice dimensions $\times$ périodes~: trajectoire de l'occultation énergétique.}\label{tab:synoptique-ch1}
\centering\scriptsize
\renewcommand{\arraystretch}{1.12}
\setlength{\tabcolsep}{3pt}
\begin{tabular}{@{}llcccccccccccc@{}}
\toprule
 & \textbf{Période} & \rotatebox{70}{\textbf{D1}} & \rotatebox{70}{\textbf{D2}} & \rotatebox{70}{\textbf{D3}} & \rotatebox{70}{\textbf{D4}} & \rotatebox{70}{\textbf{D5}} & \rotatebox{70}{\textbf{D6}} & \rotatebox{70}{\textbf{D7}} & \rotatebox{70}{\textbf{D8}} & \rotatebox{70}{\textbf{D9}} & \rotatebox{70}{\textbf{D10}} & \rotatebox{70}{\textbf{0-TER}} & \rotatebox{70}{\textbf{0-GOV}} \\
\midrule
§\ref{sec:era0} & Avant 1794 & $\bullet$ & --- & $\circ$ & $\bullet$ & $\circ$ & --- & $\circ$ & --- & $\circ$ & $\bullet$ & --- & $\circ$ \\
§\ref{sec:telegraph} & 1794--1880 & $\bullet$ & $\bullet$ & $\circ$ & $\bullet$ & $\bullet$ & $\bullet$ & $\circ$ & $\bullet$ & $\bullet$ & $\bullet$ & --- & $\bullet$ \\
§\ref{sec:calculating} & 1820--1914 & $\circ$ & $\circ$ & $\bullet$ & $\bullet$ & $\circ$ & --- & --- & $\circ$ & --- & $\circ$ & --- & --- \\
§\ref{sec:mechanography} & 1890--1940 & $\circ$ & $\bullet$ & $\bullet$ & $\bullet$ & $\circ$ & --- & $\circ$ & $\bullet$ & --- & $\bullet$ & --- & $\circ$ \\
§\ref{sec:telephone} & 1876--1920 & $\circ$ & $\bullet$ & $\circ$ & $\bullet$ & $\bullet$ & $\circ$ & $\bullet$ & $\circ$ & $\circ$ & $\bullet$ & --- & $\bullet$ \\
§\ref{sec:wireless} & 1895--1930 & $\circ$ & $\bullet$ & --- & $\circ$ & $\bullet$ & $\bullet$ & --- & $\circ$ & --- & $\bullet$ & --- & $\bullet$ \\
§\ref{sec:electrification} & 1880--1930 & --- & $\bullet$ & $\bullet$ & $\circ$ & $\bullet$ & --- & $\bullet$ & $\bullet$ & --- & $\bullet$ & $\nearrow$ & $\bullet$ \\
§\ref{sec:computing} & 1945--1975 & --- & $\bullet$ & $\bullet$ & $\circ$ & $\bullet$ & --- & --- & $\bullet$ & --- & $\bullet$ & $\circ$ & $\bullet$ \\
§\ref{sec:convergence} & 1975--2007 & $\circ$ & $\bullet$ & $\bullet$ & $\circ$ & $\bullet$ & $\bullet$ & $\bullet$ & $\bullet$ & $\circ$ & $\bullet$ & $\bullet$ & $\circ$ \\
§\ref{sec:contemporain} & 2007-- & $\circ$ & $\bullet$ & $\bullet$ & $\nearrow$ & $\bullet$ & $\circ$ & $\circ$ & $\bullet$ & $\bullet$ & $\bullet$ & $\bullet$ & $\bullet$ \\
\midrule
\multicolumn{2}{@{}l}{\textbf{Synthèse}} & \multicolumn{10}{c}{\emph{Dimensions stables~: D2, D5, D8, D10}} & $\boldsymbol{\nearrow\nearrow}$ & \\
\bottomrule
\end{tabular}
\renewcommand{\arraystretch}{1.0}
\par\medskip
\footnotesize
\textbf{Légende}~: $\bullet$~=~dimension centrale dans cette période~; $\circ$~=~dimension présente mais secondaire~; ---~=~dimension non pertinente ou non documentée~; $\nearrow$~=~dimension en émergence.
\par\smallskip
\textbf{Lecture de la colonne 0-TER}~: La dimension énergie-matière passe de «~---~» (non pertinente) dans les sections §\ref{sec:era0}--§\ref{sec:wireless}, à «~$\nearrow$~» (émergente) en §\ref{sec:electrification} avec l'électrification du substrat, à «~$\circ$~» (présente mais secondaire) en §\ref{sec:computing} où l'énergie est une contrainte d'ingénierie mais pas un coût économique visible, à «~$\bullet$~» (centrale) à partir de §\ref{sec:convergence}. Cette trajectoire illustre l'\emph{occultation progressive puis la réémergence} de la contrainte énergétique dans l'analyse des techniques de l'information.
\end{table}



% #############################################################################
%
%                    ENTRÉES BIBLIOGRAPHIQUES À AJOUTER
%                         (si absentes de Bibliography.bib)
%
% #############################################################################

% À ajouter dans Bibliography.bib si ces entrées n'existent pas déjà :

% @book{miller_chip_2022,
%   author = {Miller, Chris},
%   title = {Chip War: The Fight for the World's Most Critical Technology},
%   publisher = {Scribner},
%   year = {2022},
%   address = {New York}
% }
%
% @article{irwin_learning_1994,
%   author = {Irwin, Douglas A. and Klenow, Peter J.},
%   title = {Learning-by-Doing Spillovers in the Semiconductor Industry},
%   journal = {Journal of Political Economy},
%   volume = {102},
%   number = {6},
%   pages = {1200--1227},
%   year = {1994}
% }
%
% @book{temin_fall_1987,
%   author = {Temin, Peter and Galambos, Louis},
%   title = {The Fall of the Bell System: A Study in Prices and Politics},
%   publisher = {Cambridge University Press},
%   year = {1987},
%   address = {Cambridge}
% }
%
% @book{campbell-kelly_software_2003,
%   author = {Campbell-Kelly, Martin},
%   title = {From Airline Reservations to Sonic the Hedgehog: A History of the Software Industry},
%   publisher = {MIT Press},
%   year = {2003},
%   address = {Cambridge, MA}
% }



% #############################################################################
%
%                         RÉCAPITULATIF DES MODIFICATIONS
%
% #############################################################################
%
% ┌─────────────────────────────────────────────────────────────────────────────┐
% │ SECTION    │ FICHIER                          │ ACTION                      │
% ├─────────────────────────────────────────────────────────────────────────────┤
% │ §1.7       │ ch1_s07_electrification.tex      │                             │
% │            │   Ligne 31+                      │ + BLOC §1.7-A (Berlin/Londres)
% │            │   Lignes 33-43                   │ SUPPRIMER (commentaires)    │
% │            │   Après §Enigma (~L97)           │ + BLOC §1.7-B (Bilan T→C→S) │
% │            │   Après §1.7-B                   │ + BLOC §1.7-C (Tableau dim.) │
% ├─────────────────────────────────────────────────────────────────────────────┤
% │ §1.8       │ ch1_s08_calcul_programmable.tex  │                             │
% │            │   Après L30 (intro)              │ + BLOC §1.8-A (Transistor)  │
% │            │   Après §moore_law (~L89)        │ + BLOC §1.8-B (Consent dec.)│
% │            │   Dans bilan                     │ + BLOC §1.8-C (Tableau eff.)│
% │            │   Après §1.8-B                   │ + BLOC §1.8-D (Bilan)       │
% │            │   Après §1.8-D                   │ + BLOC §1.8-E (Tableau dim.)│
% ├─────────────────────────────────────────────────────────────────────────────┤
% │ §1.9       │ ch1_s09_convergence.tex          │                             │
% │            │   Après L99                      │ + BLOC §1.9-A (Paradoxe)    │
% │            │   Après L138 (fin)               │ + BLOC §1.9-B (Tableau dim.)│
% ├─────────────────────────────────────────────────────────────────────────────┤
% │ §1.10      │ ch1_s10_information_consomme.tex │                             │
% │            │   Ligne 52                       │ REMPLACER \wip{} (§1.10-A)  │
% │            │   Lignes 93,98,113,116,122,125,152│ REMPLACER {?} → {+} (§1.10-B)
% │            │   Après L167                     │ + BLOC §1.10-C (Tableau dim.)
% ├─────────────────────────────────────────────────────────────────────────────┤
% │ §1.11      │ ch1_s11_archetypes.tex           │                             │
% │            │   Après L83                      │ + BLOC §1.11-A (Tab. synopt.)
% ├─────────────────────────────────────────────────────────────────────────────┤
% │ Biblio     │ Bibliography.bib                 │ + 4 entrées si absentes     │
% └─────────────────────────────────────────────────────────────────────────────┘
%
% VOLUME TOTAL ESTIMÉ : ~10-11 pages de contenu nouveau
%
% =============================================================================
